\documentclass[11pt]{article}

% --- Packages ---
\usepackage[margin=1in]{geometry}
\usepackage{amsmath,amsthm,amssymb,mathtools}
\usepackage{bm}
\usepackage{enumitem}
\usepackage{microtype}
\usepackage{hyperref}
\usepackage[capitalise]{cleveref}
\usepackage[T1]{fontenc}
\usepackage{lmodern}
\usepackage{xcolor}
\usepackage{listings}
\usepackage{graphicx}
\usepackage{booktabs}
\usepackage{float}

% --- Hyperref setup ---
\hypersetup{
  colorlinks=true,
  linkcolor=blue,
  citecolor=teal,
  urlcolor=magenta,
  pdfauthor={IFMD},
  pdftitle={Multiplicity in PGF: The Structural M Operator and the Certified Contraction Constant $\Lambda_{\mathrm{rec}}$},
}

% --- Listings (Python) ---
\definecolor{codegray}{rgb}{0.5,0.5,0.5}
\definecolor{codered}{rgb}{0.7,0.1,0.1}
\definecolor{codeblue}{rgb}{0.12,0.12,0.7}
\lstdefinestyle{py}{
  language=Python,
  basicstyle=\ttfamily\small,
  keywordstyle=\color{codeblue}\bfseries,
  commentstyle=\color{codegray},
  stringstyle=\color{codered},
  showstringspaces=false,
  frame=single,
  framerule=0.3pt,
  breaklines=true,
  columns=fullflexible
}

% --- Macros (Notation Standardization) ---
% Multiplicity Framework Macros
\newcommand{\Sig}{\mathrm{Sig}}
\newcommand{\Mfun}{M_{\mathrm{fun}}}
\newcommand{\Mint}{M_{\mathrm{int}}}
\newcommand{\Mhat}{\widehat{M}}
\newcommand{\Mdiag}{\mathbf{M}}
\newcommand{\Lrec}{\Lambda_{\mathrm{rec}}}
\newcommand{\Lglobal}{\Lambda_{\mathrm{global}}}
\newcommand{\Llocal}{\Lambda_{\mathrm{local}}}
\newcommand{\Lcrit}{\Lambda_{\mathrm{crit}}}
\newcommand{\Lhyb}{\Lambda_{\mathrm{rec}}^{\mathrm{hyb}}}
% Prime-Graded Framework
\newcommand{\PGalgebra}{\mathcal{A}}
\newcommand{\PGprojector}{P}
\newcommand{\PGstate}{\varphi}
\newcommand{\PGspace}{\mathcal{H}}
% Utility
\newcommand{\norm}[1]{\left\lVert #1\right\rVert}
\newcommand{\abs}[1]{\left\lvert #1\right\rvert}
\newcommand{\ip}[2]{\left\langle #1,\,#2\right\rangle}

% --- Theorem environments ---
\newtheorem{theorem}{Theorem}
\newtheorem{lemma}[theorem]{Lemma}
\newtheorem{proposition}[theorem]{Proposition}
\newtheorem{corollary}[theorem]{Corollary}
\theoremstyle{definition}
\newtheorem{definition}[theorem]{Definition}
\newtheorem{assumption}[theorem]{Assumption}
\newtheorem{remark}[theorem]{Remark}
\begin{document}
\begin{titlepage}
    \centering

    {\Huge\bfseries Multiplicity in PGF \par}
    \vspace{0.5cm}
    {\LARGE\bfseries The Structural $M$ Operator and the Certified Contraction Constant $\Lrec$ \par}

  \vspace{2cm}
    {\LARGE\itshape Ryan O. Van Gelder \par}
    \vspace{1cm}
    {\Large\bfseries Citizen Gardens} \\
    \vspace{0.5cm}
    \Large{Institute for Mathematical Discovery} \\ 
    \vspace{0.5cm}
    {\normalsize \today \par}

 \vfill

    {\bfseries Abstract \par}
    \vspace{0.5em}
    \begin{minipage}{0.85\textwidth}
        \small
We present a comprehensive, Prime-Graded Framework (PGF) for multiplicity in recursive systems that strictly separates the \emph{structural} Multiplicity operator layer ($\Mfun$, $\Mint$, $\Mhat$, $\Mdiag$) from the \emph{analytic} contraction layer governed by the multiplicity (contraction) constant $\Lrec$. We define the critical boundary $\Lcrit=1/\rho(S)$, provide a certified global choice $\Lglobal(\gamma)=\gamma/\norm{S}$ with $0<\gamma<1$, introduce a Rayleigh-based local gate $\Llocal(\gamma;T)=\gamma/r(T)$, and combine them in a safe hybrid policy $\Lhyb=\min\{\Lglobal,\Llocal\}$. We prove stability and error bounds, give boundedness criteria for infinite prime families, and provide production-ready Python reference implementations.

 \end{minipage}
\end{titlepage}
\newpage

\tableofcontents


\section{Introduction \& Motivation}
% ===============================
Complex prime-weighted recursive systems appear in mathematics, physics-inspired models, and structured AI pipelines. Prior drafts mixed structural multiplicity (prime accounting, typing) with analytic gains controlling contraction, creating type-mismatches and unprovable claims. This article consolidates the theory into a PGF-aligned document with:
\begin{itemize}[leftmargin=*,nosep]
  \item a \emph{structural} layer: the $M$-objects $\Mfun,\Mint,\Mhat,\Mdiag$ that certify prime structure,
  \item an \emph{analytic} layer: the scalar $\Lrec$ that ensures contraction and stability.
\end{itemize}
Our main contributions are: (i) precise definitions and proofs; (ii) certified computation algorithms; (iii) minimal, reproducible case studies.

% ===============================
\section{Mathematical Preliminaries}
% ===============================
\subsection{Prime-Graded Framework (PGF)}
\begin{definition}[Signature Category]
Let $\Sig=\{e:\mathbb{P}\to\mathbb{Z}\ \text{with finite support}\}$ equipped with
\begin{itemize}[leftmargin=*,nosep]
  \item Tensor: $(e\otimes f)(p)=e(p)+f(p)$,
  \item Unit: $0$,
  \item Dual: $e^\vee=-e$.
\end{itemize}
\end{definition}

\begin{definition}[Prime-Graded $*$-Algebra]
A $*$-algebra $\PGalgebra=\bigoplus_{p\in\mathbb{P}}\PGalgebra_p$ with bounded projectors $\PGprojector_p:\PGalgebra\to\PGalgebra_p$, $\PGprojector_p^2=\PGprojector_p$, and a positive linear functional $\PGstate:\PGalgebra\to\mathbb{C}$.
\end{definition}

\paragraph{Hilbert structure for Rayleigh.}
We endow $\PGalgebra$ with the inner product $\ip{X}{Y}:=\PGstate(X^\*Y)$ and norm $\norm{X}^2=\PGstate(X^\*X)$. This makes the Rayleigh average $r(T)=\ip{T}{ST}/\norm{T}^2$ well-defined whenever $S$ is bounded.

\subsection{Recursion Operator and Dynamics}
Let $P_N\subset\mathbb{P}$ be a finite active set (or assume boundedness otherwise). Define
\begin{equation}
  S:=\sum_{p\in P_N} p^{\alpha_p}\PGprojector_p.
\end{equation}
The affine recursion is
\begin{equation}\label{eq:rec}
  T_{t+1}=\Lrec\,S\,T_t + F.
\end{equation}

\begin{assumption}[Settings]\label{ass:settings}
We work on a Banach space with submultiplicative norm and its induced operator norm (Hilbert structure when using Rayleigh quotients). Projectors satisfy $\norm{\PGprojector_p}=1$. For orthogonal resolution we additionally assume $\PGprojector_p\PGprojector_q=\delta_{pq}\PGprojector_p$ and $\sum_{p\in P_N}\PGprojector_p=I$.
\end{assumption}

% ===============================
\section{The Multiplicity Operator $M$ (Structural Layer)}
% ===============================
\subsection{Four Consistent Faces}
\begin{definition}[Functorial $M$]\label{def:Mfun}
$\Mfun:\Sig\to\mathbb{Q}^\times$ with $\Mfun(e)=\prod_{p}p^{e(p)}$.
\end{definition}
\begin{theorem}[Monoidal Properties]\label{thm:monoidal}
$\Mfun$ is strict monoidal: $\Mfun(0)=1$, $\Mfun(e\otimes f)=\Mfun(e)\Mfun(f)$, and $\Mfun(e^\vee)=\Mfun(e)^{-1}$.
\end{theorem}

\begin{definition}[Integer Ledger]
$\Mint(X,p)=e(p)\in\mathbb{Z}$ tracks exact prime exponents in signatures.
\end{definition}

\begin{definition}[Sector Weights]\label{def:Mhat}
For $T\in\PGalgebra$, set $\mu_p(T):=\PGstate\big((\PGprojector_pT)^\*(\PGprojector_pT)\big)$ and
\[
  \Mhat(T,p)=\frac{\mu_p(T)}{\sum_{q\in P_N}\mu_q(T)}\in[0,1],\qquad \sum_{p\in P_N}\Mhat(T,p)=1.
\]
\end{definition}

\begin{definition}[Diagonal Signature Gauge]
On a signature Hilbert space $\PGspace$, $\Mdiag\ket{e}=\Mfun(e)\ket{e}$.
\end{definition}

\begin{remark}[Separation of concerns]
$\Mfun,\Mint,\Mhat,\Mdiag$ are \emph{structural} and must \emph{never} be multiplied by real analytic contraction gains; only $\Lrec$ scales the recursion.
\end{remark}

% ===============================
\section{The Contraction Constant $\Lrec$ (Analytic Layer)}
% ===============================
\subsection{Critical and Certified Global Scales}
\begin{definition}[Critical boundary and certified global]\label{def:lambdas}
Let $S=\sum_{p}p^{\alpha_p}\PGprojector_p$. Define
\[
  \Lcrit:=\frac{1}{\rho(S)},\qquad
  \boxed{\ \Lglobal(\gamma):=\frac{\gamma}{\norm{S}}\ },\quad 0<\gamma<1.
\]
Then $\norm{\Lglobal(\gamma)S}\le\gamma<1$.
\end{definition}

\begin{remark}[Diagonal fast-path]\label{rem:diag}
Under orthogonal resolution, $S$ is self-adjoint positive with spectrum $\sigma(S)=\{p^{\alpha_p}\}_p$ and $\norm{S}=\rho(S)=\max_p p^{\alpha_p}$. Hence $\Lglobal(\gamma)=\gamma/\max_p p^{\alpha_p}$ and $\norm{\Lglobal(\gamma)S}=\gamma$.
\end{remark}

\subsection{Local Rayleigh Gate and Hybrid Policy}
\begin{definition}[Rayleigh average]
$r(T):=\dfrac{\ip{T}{ST}}{\norm{T}^2}=\sum_{p\in P_N}\Mhat(T,p)\,p^{\alpha_p}$ (orthogonal case).
\end{definition}
\begin{definition}[Local and hybrid policies]
For $\gamma\in(0,1)$ and $r(T)>0$, set
\[
  \boxed{\ \Llocal(\gamma;T):=\frac{\gamma}{r(T)}\ },\qquad
  \boxed{\ \Lhyb(T):=\min\{\Lglobal(\gamma),\,\Llocal(\gamma;T)\}\ }.
\]
If $S$ is non-normal, compute $r(T)$ using $S_{\mathrm{sym}}=(S+S^\*)/2$ and keep $\Lglobal$ as a hard fence.
\end{definition}

\subsection{Stability and Error Bounds}
\begin{theorem}[Global certified stability]\label{thm:global}
With $0<\gamma<1$ and $\Lrec=\Lglobal(\gamma)$, the map $\Phi(T)=\Lrec S T + F$ is a contraction with constant $c=\norm{\Lrec S}\le\gamma<1$. Thus $T_\infty=(I-\Lrec S)^{-1}F$ is unique and
\[
\norm{T_t-T_\infty}\le \gamma^t\norm{T_0-T_\infty},\qquad
\norm{T_\infty}\le \frac{\norm{F}}{1-\gamma}.
\]
\end{theorem}
\begin{proof}
Immediate from submultiplicativity: $\norm{\Phi(T)-\Phi(T')}=\norm{\Lrec S(T-T')}\le\gamma\norm{T-T'}$.
\end{proof}

\begin{lemma}[Directional local gate]\label{lem:local}
For self-adjoint $S$, with $\Llocal(\gamma;T)=\gamma/r(T)$ one has $\norm{(\Llocal S)T}=\gamma\norm{T}$.
\end{lemma}

\begin{proposition}[Hybrid safety]\label{prop:hybrid}
$\Lhyb(T)=\min\{\Lglobal(\gamma),\Llocal(\gamma;T)\}$ satisfies $\norm{\Lhyb(T)S}\le\gamma<1$ and never exceeds the global fence.
\end{proposition}

\subsection{Infinite Prime Families: Boundedness Criterion}
\begin{proposition}[Boundedness for $|P_N|=\infty$]\label{prop:infinite}
Assume orthogonal resolution with $\sum_{p\in P_N}\PGprojector_p=I$ and $S=\sum_p p^{\alpha_p}\PGprojector_p$. Then $S$ is bounded iff $\sup_{p}p^{\alpha_p}<\infty$. For unbounded primes this is equivalent to $\limsup_{p\to\infty}\alpha_p\le 0$, and then $\norm{S}=\sup_p p^{\alpha_p}$.
\end{proposition}
\begin{proof}
$S$ is diagonal in the $\PGprojector_p$ basis with eigenvalues $p^{\alpha_p}$, so $\norm{S}=\sup_p p^{\alpha_p}$.
\end{proof}

\begin{remark}[Golden-ratio anecdote]
If at a fixed point $T_\infty$ the Rayleigh average satisfies $r(T_\infty)\approx\varphi^{-1}$, then the \emph{largest certified local} value is $\Llocal^\star(\gamma;T_\infty)=\gamma\,\varphi$. This is contingent on spectra and sector weights, not universal.
\end{remark}

% ===============================
\section{Certified Computation: Reference Implementations}
% ===============================
We provide minimal, production-ready snippets. Projectors are linear, idempotent diagonal masks; $\PGstate$ is quadratic, ensuring the Hilbert structure.

\subsection{Framework and Sector Weights}
\begin{lstlisting}[style=py,caption={PrimeGradedFramework with linear idempotent projectors and quadratic state},label={lst:framework}]
import numpy as np

class PrimeGradedFramework:
    def __init__(self, primes, masks):
        """
        primes: list of primes
        masks: dict prime -> boolean/float mask of shape (d,)
        """
        self.primes = primes
        self.masks = masks
        self.dim = len(next(iter(masks.values())))
        # Linear idempotent projectors: P_p(x) = mask * x
        self.projectors = {p: (lambda m: (lambda x: m * x))(masks[p]) for p in primes}

    # φ(X) = ||X||_2^2 ensures <X,Y> = φ(X*Y) for vectors
    def phi(self, x):
        return float(np.vdot(x, x))

    def sector_weight(self, T, p):
        Tp = self.projectors[p](T)
        mu_p = self.phi(Tp)
        mu_tot = sum(self.phi(self.projectors[q](T)) for q in self.primes)
        return (mu_p / mu_tot) if mu_tot > 0 else 0.0

    def zero_element(self):
        return np.zeros(self.dim, dtype=float)

    def zero_element_like(self, x):
        return np.zeros_like(x)
\end{lstlisting}

\subsection{Global Certificate (Diagonal Fast-Path) and Safe Norm Upper Bound}
\begin{lstlisting}[style=py,caption={Diagonal $\|S\|$ and one-sided upper bound for general $S$},label={lst:norms}]
def lambda_global_diag(primes, alpha, gamma=0.9):
    smax = max(p ** alpha.get(p, 0.0) for p in primes)
    if smax <= 0:
        raise ValueError("Diagonal S requires positive weights.")
    return gamma / smax

def upper_bound_norm(apply_S, dim, iters=50, probes=8, seed=0):
    """
    One-sided upper bound for ||S||_2 (never underestimates).
    """
    rng = np.random.default_rng(seed)
    ub = 0.0
    for _ in range(probes):
        x = rng.standard_normal(dim); x /= np.linalg.norm(x)
        for _ in range(iters):
            y = apply_S(x); n = np.linalg.norm(y)
            if n == 0: break
            x = y / max(n, 1e-300)
        ub = max(ub, np.linalg.norm(apply_S(x)))
    return max(ub, 1e-12)
\end{lstlisting}

\subsection{Local Gate, Hybrid Policy, and Certified Update}
\begin{lstlisting}[style=py,caption={Local Rayleigh gate, hybrid selection, and certified update},label={lst:hybrid}]
def lambda_local_gamma(T, primes, alpha, framework, gamma=0.9, eps=1e-12):
    r = sum(framework.sector_weight(T, p) * (p ** alpha.get(p, 0.0)) for p in primes)
    if not np.isfinite(r) or r <= eps:
        return None
    return gamma / r

def lambda_hybrid(T, lambda_global, local_fn=None, hard_cap=None):
    vals = [lambda_global]
    if local_fn is not None:
        lc = local_fn(T)
        if lc is not None:
            vals.append(lc)
    v = min(vals)
    return min(v, hard_cap) if hard_cap is not None else v

def certified_update(T, F, primes, alpha, projectors, framework, Lambda_global,
                     local_fn=None, hard_cap=None):
    Lambda = lambda_hybrid(T, Lambda_global, local_fn=local_fn, hard_cap=hard_cap)
    T_next = framework.zero_element_like(T)
    for p in primes:
        T_next += Lambda * (p ** alpha.get(p, 0.0)) * projectors[p](T)
    T_next += F
    return T_next
\end{lstlisting}

\subsection{End-to-End: Diagonal Case Study (Reproducible)}
\begin{lstlisting}[style=py,caption={Diagonal system: convergence with $\|\Lrec S\|\le\gamma$},label={lst:diagdemo}]
def demo_diagonal_convergence():
    primes = [2, 3, 5]
    d = 12
    masks = {
        2: np.array([1]*4 + [0]*8, dtype=float),
        3: np.array([0]*4 + [1]*4 + [0]*4, dtype=float),
        5: np.array([0]*8 + [1]*4, dtype=float),
    }
    fw = PrimeGradedFramework(primes, masks)
    alpha = {2: -1.0, 3: -1.0, 5: -1.0}  # S diag entries: {1/2, 1/3, 1/5}
    gamma = 0.8

    # Diagonal fast-path: ||S|| = max_p p^{alpha_p} = 1/2
    Lglob = lambda_global_diag(primes, alpha, gamma=gamma)  # 0.8 / 0.5 = 1.6
    local_fn = lambda T: lambda_local_gamma(T, primes, alpha, fw, gamma=gamma)

    T = np.linspace(1.0, 2.0, fw.dim)
    F = 0.1 * np.ones(fw.dim)

    rates = []
    for _ in range(50):
        T_new = certified_update(T, F, primes, alpha, fw.projectors, fw, Lglob, local_fn=local_fn)
        rates.append(np.linalg.norm(T_new - T))
        T = T_new

    print("Λ_global =", Lglob, " (Λ_global>1 is fine since ||S||<1 ⇒ ||ΛS||≤γ)")
    print("final ||ΔT||:", rates[-1])
\end{lstlisting}

\subsection{Sketch: Non-Diagonal Case (Safe Norm Fence)}
\begin{lstlisting}[style=py,caption={Building apply_S and fencing with a one-sided upper bound},label={lst:nondiag}]
def build_apply_S(primes, alpha, projectors, framework):
    def apply_S(x):
        out = framework.zero_element_like(x)
        for p in primes:
            out += (p ** alpha.get(p, 0.0)) * projectors[p](x)
        return out
    return apply_S

def demo_nondiagonal_certified(primes, alpha, projectors, framework, gamma=0.8):
    apply_S = build_apply_S(primes, alpha, projectors, framework)
    ub = upper_bound_norm(apply_S, dim=framework.dim, iters=40, probes=8)
    Lglob = gamma / ub
    local_fn = lambda T: lambda_local_gamma(T, primes, alpha, framework, gamma=gamma)
    # proceed as in the diagonal demo...
\end{lstlisting}

% ===============================
\section{Applications}
% ===============================
\subsection{Case Study: Diagonal System}
\textbf{Parameters:} $P_3=\{2,3,5\}$, $\alpha_2=\alpha_3=\alpha_5=-1$, $\gamma=0.8$, $F=0.1\cdot \mathbf{1}$, disjoint masks of size $4$ each.

\textbf{Certificates:} $\norm{S}=\max\{2^{-1},3^{-1},5^{-1}\}=1/2$, hence $\Lglobal=0.8/(1/2)=1.6$ and $\norm{\Lglobal S}\le 0.8$.

\textbf{Outcome:} Empirical linear rate $\approx \gamma=0.8$; convergence to $T_\infty=(I-\Lglobal S)^{-1}F$; $\Llocal$ further tightens steps statewise without breaching the global fence.

\subsection{Case Study: Non-Diagonal System}
\textbf{Challenge:} Non-orthogonal projectors (e.g., overlapping masks or learned linear maps) break exact diagonalization.

\textbf{Solution:} Use \texttt{upper\_bound\_norm} to fence $\norm{S}$ from above, set $\Lglobal=\gamma/\text{UB}$, and optionally deploy $S_{\mathrm{sym}}$ in the local Rayleigh average. Hybrid policy preserves $\norm{\Lrec S}\le\gamma$.

% ===============================
\section{Conclusion and Outlook}
% ===============================
We completed a PGF-consistent split between the structural $M$ layer and the analytic contraction scalar $\Lrec$. This yields clean, certifiable stability results, a minimal set of assumptions, and simple, robust algorithms. Future directions include learned local gates, multi-scale $\Lrec$ spectra, and domain-specific PETC constraints.

% ===============================
\appendix
\section*{Appendix A: Quick Reference (Definitions \& Results)}
\begin{itemize}[leftmargin=*,nosep]
  \item $S=\sum_p p^{\alpha_p}\PGprojector_p$; recursion $T_{t+1}=\Lrec ST_t+F$.
  \item $\Lcrit=1/\rho(S)$ (boundary), $\Lglobal(\gamma)=\gamma/\norm{S}$ (certificate).
  \item $r(T)=\ip{T}{ST}/\norm{T}^2=\sum_p \Mhat(T,p)p^{\alpha_p}$ (orthogonal case).
  \item $\Llocal(\gamma;T)=\gamma/r(T)$, $\Lhyb=\min\{\Lglobal,\Llocal\}$.
  \item If orthogonal: $\norm{S}=\max_p p^{\alpha_p}$.
  \item Infinite primes: $S$ bounded $\Leftrightarrow \sup_p p^{\alpha_p}<\infty$.
\end{itemize}

\section*{Appendix B: Minimal Make Tips}
\begin{itemize}[leftmargin=*,nosep]
  \item Compile with \texttt{pdflatex} $\times$2 (for cross-refs).
  \item Use \texttt{\textbackslash lstinputlisting} to include full Python files if desired.
\end{itemize}

\nocite{*}
\bibliographystyle{unsrt}
\bibliography{references}
\end{document}




